\uuid{toqE}
\titre{Nombres premiers}
\auteur{bodin ; croizat ; sacre}
\organisation{exo7}
\licence{CC BY-NC-SA 4.0}
\datecreate{2022-03-08}
\niveau{L1}
\module{Algèbre}
\chapitre{Arithmétique dans Z}
\sousChapitre{Nombres premiers, nombres premiers entre eux}
\difficulte{2}
\theme{Nombres premiers}
\exo7qcmid{hilisit.arithmetique.030}

\contenu{
\begin{question}
\qclassification{103.04}
    Quelles sont les affirmations vraies ?
    \begin{answers}
        \good{La somme de deux nombres premiers $\ge3$ n'est jamais un nombre premier.}
        \good{Le produit de deux nombres premiers $\ge3$ n'est jamais un nombre premier.}
        \bad{Il existe un nombre premier $p\ge3$ tel que $p+1$ soit aussi premier.}
        \good{Il existe un nombre premier $p\ge3$ tel que $p+2$ soit aussi premier.}
    \end{answers}
    \begin{explanations}
        Le produit de deux nombres premiers n'est jamais un nombre premier (par définition de ce qu'est un nombre premier).
        Pour la somme, cela peut arriver, par exemple $2+3=5$, mais pour deux nombres premiers $\ge3$, ils sont impairs, donc la somme est paire et n'est pas un nombre premier.
        De même si $p\ge3$ est premier, il est impair, donc $p+1$ est pair et n'est pas premier.
        Par contre pour $p=11$ alors $p+2=13$ est aussi premier, d'autres exemples sont $17$ et $19$ ou bien $101$ et $103$.
    \end{explanations}
\end{question}
}
